% RTT Presentation of Y_hbar(sl_3) from the Ordered Bar Complex d^2=0
% Computation record: explicit derivation from first principles
% Generated by sl3_rtt_derivation.py (verified numerically)

\subsection*{RTT presentation of $Y_\hbar(\mathfrak{sl}_3)$
 from the ordered bar complex $d^2=0$}

\subsubsection*{1.\ The Casimir tensor}

The Cartan matrix of $\mathfrak{sl}_3$ is
$A = \bigl(\begin{smallmatrix} 2 & -1 \\ -1 & 2 \end{smallmatrix}\bigr)$
with inverse
$A^{-1} = \tfrac{1}{3}\bigl(\begin{smallmatrix} 2 & 1 \\ 1 & 2 \end{smallmatrix}\bigr)$.
The split Casimir tensor is
\begin{equation}\label{eq:sl3-casimir}
\Omega \;=\;
 \sum_{\alpha > 0} (e_\alpha \otimes f_\alpha + f_\alpha \otimes e_\alpha)
 \;+\; \sum_{i,j=1}^{2} (A^{-1})_{ij}\, h_i \otimes h_j.
\end{equation}
Explicitly:
\begin{align*}
\Omega &= e_1 \otimes f_1 + f_1 \otimes e_1
 + e_2 \otimes f_2 + f_2 \otimes e_2
 + e_{12} \otimes f_{12} + f_{12} \otimes e_{12} \\
 &\quad + \tfrac{2}{3}\, h_1 \otimes h_1
 + \tfrac{1}{3}(h_1 \otimes h_2 + h_2 \otimes h_1)
 + \tfrac{2}{3}\, h_2 \otimes h_2.
\end{align*}
In the defining representation $V = \mathbb{C}^3$,
using $\mathfrak{gl}_3$ matrix units
$e_1 = E_{12}$, $e_2 = E_{23}$, $e_{12} = E_{13}$,
$f_1 = E_{21}$, $f_2 = E_{32}$, $f_{12} = E_{31}$,
$h_1 = E_{11} - E_{22}$, $h_2 = E_{22} - E_{33}$:
\begin{equation}\label{eq:casimir-perm}
\Omega \;=\; P - \tfrac{1}{3}\, I_9,
\end{equation}
where $P \colon V \otimes V \to V \otimes V$ is the permutation
$P|i,j\rangle = |j,i\rangle$ and $I_9$ is the identity on
$V \otimes V$. The eigenvalues of $\Omega$ are
$\tfrac{2}{3}$ (multiplicity~$6$, symmetric part)
and $-\tfrac{4}{3}$ (multiplicity~$3$, antisymmetric part),
with $\operatorname{Tr}(\Omega) = 0$.
Verified numerically: $\|\Omega - (P - I_9/3)\|_{\max} < 10^{-15}$.

\subsubsection*{2.\ The collision residue $r$-matrix}

For $\widehat{\mathfrak{sl}}_3$ at level~$k$, the OPE
$J^a(z)\,J^b(w) \sim k\,\kappa^{ab}/(z-w)^2 + f^{ab}{}_c\,J^c(w)/(z-w)$
has maximum pole order~$2$.
The bar propagator $d\log(z_1 - z_2)$ absorbs one power:
\[
 z^{-2} \;\mapsto\; z^{-1}, \qquad
 z^{-1} \;\mapsto\; z^{0}\;\text{(regular)}.
\]
The collision residue is therefore
\begin{equation}\label{eq:sl3-r-matrix}
r(z) \;=\; \frac{k\,\Omega}{z}
\;=\; \frac{k}{z}\Bigl(P - \tfrac{1}{3}\,I_9\Bigr),
\end{equation}
a single simple pole. This satisfies the classical Yang--Baxter equation:
the infinitesimal braid relation
$[\Omega_{12}, \Omega_{13} + \Omega_{23}] = 0$
holds by ad-invariance of the Killing form (verified numerically).

\subsubsection*{3.\ Yang's $R$-matrix (9$\times$9)}

The quantization of $r(z) = k\Omega/z$ is the Yang $R$-matrix
\begin{equation}\label{eq:yang-r-sl3}
R(u) \;=\; u\,I_9 + \hbar\,P_9,
\end{equation}
acting on $\mathbb{C}^3 \otimes \mathbb{C}^3$ with $\hbar = 1/(k + 3)$
(shifted by the dual Coxeter number $h^\vee = 3$).
Component form:
\begin{equation}\label{eq:R-components}
R(u)_{(ij),(kl)}
\;=\; u\,\delta_{ik}\,\delta_{jl}
\;+\; \hbar\,\delta_{il}\,\delta_{jk}.
\end{equation}
The Yang--Baxter equation
$R_{12}(u{-}v)\,R_{13}(u)\,R_{23}(v)
= R_{23}(v)\,R_{13}(u)\,R_{12}(u{-}v)$
on $(\mathbb{C}^3)^{\otimes 3} = \mathbb{C}^{27}$
is verified numerically at all tested spectral values.

\subsubsection*{4.\ The RTT relation}

The transfer matrix
$T(u) = (t_{ij}(u))_{1 \le i,j \le 3}
\in \operatorname{End}(\mathbb{C}^3) \otimes Y_\hbar(\mathfrak{gl}_3)[\![u^{-1}]\!]$
with $t_{ij}(u) = \delta_{ij} + \sum_{r \ge 1} t_{ij}^{(r)}\,u^{-r}$ satisfies
\begin{equation}\label{eq:rtt-sl3}
R(u{-}v)\,T_1(u)\,T_2(v) \;=\; T_2(v)\,T_1(u)\,R(u{-}v).
\end{equation}
In components, this is equivalent to
\begin{equation}\label{eq:rtt-component}
(u - v)\,[t_{ik}(u),\, t_{jl}(v)]
\;=\;
\hbar\bigl(t_{jk}(v)\,t_{il}(u) - t_{jk}(u)\,t_{il}(v)\bigr).
\end{equation}

\subsubsection*{5.\ Expanded RTT at order $(p,q)$}

Expanding $T(u) = \sum_{r \ge 0} T^{(r)} u^{-r}$ with $T^{(0)} = I_3$
and equating coefficients of $u^{-p}\,v^{-q}$:
\begin{equation}\label{eq:rtt-pq}
\bigl[t^{(p+1)}_{ik},\, t^{(q)}_{jl}\bigr]
- \bigl[t^{(p)}_{ik},\, t^{(q+1)}_{jl}\bigr]
\;=\;
\hbar\bigl(t^{(q)}_{jk}\,t^{(p)}_{il}
 - t^{(p)}_{jk}\,t^{(q)}_{il}\bigr).
\end{equation}

\noindent\textbf{At order $(1,0)$} (level-1 relation):
\begin{equation}\label{eq:level-1-sl3}
\bigl[t^{(1)}_{ik},\, t^{(1)}_{jl}\bigr]
\;=\;
\hbar\bigl(\delta_{il}\,t^{(1)}_{jk}
 - \delta_{jk}\,t^{(1)}_{il}\bigr).
\end{equation}
This is the $\mathfrak{gl}_3$ Lie bracket, confirming that
$\{t^{(1)}_{ij}\}$ form a copy of $\mathfrak{gl}_3$.

\noindent\textbf{At order $(1,1)$} (symmetry constraint):
$\bigl[t^{(2)}_{ik},\, t^{(1)}_{jl}\bigr]
= \bigl[t^{(1)}_{ik},\, t^{(2)}_{jl}\bigr]$.

\noindent\textbf{At order $(2,1)$} (cubic Serre):
$\bigl[t^{(3)}_{ik},\, t^{(1)}_{jl}\bigr]
- \bigl[t^{(2)}_{ik},\, t^{(2)}_{jl}\bigr]
= \hbar\bigl(t^{(1)}_{jk}\,t^{(2)}_{il}
 - t^{(2)}_{jk}\,t^{(1)}_{il}\bigr)$.

\subsubsection*{6.\ Relation count}

\begin{center}
\begin{tabular}{lcc}
\hline
& $\mathfrak{gl}_3$ & $\mathfrak{sl}_3$ \\
\hline
Generators per level & $9$ & $8$ \\
Raw RTT components & $9^2 = 81$ & --- \\
Independent (antisymmetry) & $\binom{9}{2} = 36$ & $\binom{8}{2} = 28$ \\
\quad Commuting pairs ($[\cdot,\cdot]=0$) & $15$ & --- \\
\quad Nontrivial brackets & $21$ & --- \\
Total generators (levels $0{+}1$) & $18$ & $16$ \\
\hline
\end{tabular}
\end{center}

\noindent
All $36$ level-1 relations for $\mathfrak{gl}_3$ are listed
explicitly in the computation log. The passage from
$\mathfrak{gl}_3$ to $\mathfrak{sl}_3$ imposes the quantum
determinant condition $\operatorname{qdet} T(u) = 1$.

\subsubsection*{7.\ Complete level-1 bracket table}

The $21$ nontrivial brackets (the remaining $15$ are commuting pairs):
\begin{align*}
[t_{11}, t_{12}] &= -\hbar\,t_{12},
&[t_{11}, t_{13}] &= -\hbar\,t_{13},
&[t_{11}, t_{21}] &= \hbar\,t_{21}, \\
[t_{11}, t_{31}] &= \hbar\,t_{31},
&[t_{12}, t_{21}] &= \hbar(t_{22} - t_{11}),
&[t_{12}, t_{22}] &= -\hbar\,t_{12}, \\
[t_{12}, t_{23}] &= -\hbar\,t_{13},
&[t_{12}, t_{31}] &= \hbar\,t_{32},
&[t_{13}, t_{21}] &= \hbar\,t_{23}, \\
[t_{13}, t_{31}] &= \hbar(t_{33} - t_{11}),
&[t_{13}, t_{32}] &= -\hbar\,t_{12},
&[t_{13}, t_{33}] &= -\hbar\,t_{13}, \\
[t_{21}, t_{22}] &= \hbar\,t_{21},
&[t_{21}, t_{32}] &= \hbar\,t_{31},
&[t_{22}, t_{23}] &= -\hbar\,t_{23}, \\
[t_{22}, t_{32}] &= \hbar\,t_{32},
&[t_{23}, t_{31}] &= -\hbar\,t_{21},
&[t_{23}, t_{32}] &= \hbar(t_{33} - t_{22}), \\
[t_{23}, t_{33}] &= -\hbar\,t_{23},
&[t_{31}, t_{33}] &= \hbar\,t_{31},
&[t_{32}, t_{33}] &= \hbar\,t_{32}.
\end{align*}
The $15$ commuting pairs include all $(t_{ij}, t_{kl})$ with
$i \ne l$ and $j \ne k$ (orthogonal matrix units).

\subsubsection*{8.\ Identification with $\mathfrak{sl}_3$ generators}

\[
\begin{array}{cccccccc}
t_{12} & t_{23} & t_{13} & t_{21} & t_{32} & t_{31}
 & t_{11}{-}t_{22} & t_{22}{-}t_{33} \\
\| & \| & \| & \| & \| & \| & \| & \| \\
e_1 & e_2 & e_{12} & f_1 & f_2 & f_{12} & h_1 & h_2
\end{array}
\]
Under this identification:
\begin{itemize}
\item $[t_{12}, t_{23}] = \hbar\,t_{13}$ $\longleftrightarrow$
 $[e_1, e_2] = e_{12}$;
\item $[t_{12}, t_{21}] = \hbar(t_{22} - t_{11})$ $\longleftrightarrow$
 $[e_1, f_1] = h_1$;
\item $[t_{12}, t_{32}] = 0$ $\longleftrightarrow$
 $[e_1, f_2] = 0$.
\end{itemize}

\subsubsection*{9.\ Triangle sectors}

The degree-$3$ ordered bar complex
$\Barch_3^{\mathrm{ord}}(\widehat{\mathfrak{sl}}_3)$
exhibits four \emph{triangle sectors}---triples of root vectors
whose roots form a closed triangle
$\alpha + \beta = \gamma$ in the root system:
\[
T^L_{12} = E_{13}^{(1)}\,E_{21}^{(2)}\,E_{32}^{(3)},
\qquad
T^L_{21} = E_{13}^{(1)}\,E_{32}^{(2)}\,E_{21}^{(3)};
\]
\[
T^R_{12} = E_{12}^{(1)}\,E_{23}^{(2)}\,E_{31}^{(3)},
\qquad
T^R_{21} = E_{23}^{(1)}\,E_{12}^{(2)}\,E_{31}^{(3)}.
\]

The triangle coefficient is
\begin{equation}\label{eq:triangle-half}
\lambda^L = \lambda^R
= \frac{c_1\,c_2}{2\,c_{12}}
= \frac{1}{2}
= \int_0^1 (1-t)\,dt
= B(1,2),
\end{equation}
the $n=2$ beta integral on the $\operatorname{FM}_3(\mathbb{R})$ ordered cell.
This arises from the parametrization of the FM cell
by the cross-ratio $t = (z_2 - z_1)/(z_3 - z_1) \in [0,1]$
for three ordered points $z_1 < z_2 < z_3$.

Explicitly, on $[e_1 | e_2 | f_{12}]$:
\begin{align*}
d_L[e_1|e_2|f_{12}] &= [e_{12} | f_{12}], \\
d_R[e_1|e_2|f_{12}] &= [e_1 | f_1], \\
d^2 &= d_R\,d_L - d_L\,d_R
= [h_1 + h_2] - [h_1] = [h_2] \neq 0
\end{align*}
before the FM weight correction. The coefficient $\tfrac{1}{2}$
from the beta integral ensures $d^2 = 0$ after integrating over
the FM cell.

\subsubsection*{10.\ Serre relations from root-space vanishing}

The Serre relation $[e_1, [e_1, e_2]] = 0$ follows from:
\[
2\alpha_1 + \alpha_2 \notin \Phi(\mathfrak{sl}_3)
\quad\implies\quad
(\mathfrak{sl}_3)_{2\alpha_1 + \alpha_2} = 0
\quad\implies\quad
\text{the degree-3 bar class } [e_1|e_1|e_2]
\text{ is exact}.
\]
Similarly $[e_2,[e_2,e_1]] = 0$ from
$\alpha_1 + 2\alpha_2 \notin \Phi$.
These are the only Serre relations for $\mathfrak{sl}_3$
(the Dynkin diagram $A_2$ has no non-adjacent nodes).

\subsubsection*{11.\ Bar complex interpretation}

The RTT relation~\eqref{eq:rtt-sl3} is equivalent to $d^2 = 0$
on $\Barch_3^{\mathrm{ord}}(\widehat{\mathfrak{sl}}_3)$:
\begin{itemize}
\item The bar differential
 $d = d_L - d_R$ has two face maps
 (left and right collision, controlled by $r(z) = k\Omega/z$).
\item $d^2 = 0$ requires $d_L d_R = d_R d_L$ on degree-$3$ elements.
\item At the classical level: this is the CYBE.
\item At the quantum level: this is the RTT relation.
\item The $T$-matrix encodes bar cohomology classes
 (the ``insertion of a bar generator in the third slot'').
\end{itemize}

\subsubsection*{Summary of results}

\begin{enumerate}
\item $\mathfrak{sl}_3$ data: $8$ generators,
 Killing form and Jacobi identity verified.
\item Casimir: $\Omega = P - \frac{1}{3} I_9$
 in the defining representation ($9 \times 9$).
\item $R$-matrix: $R(u) = u\,I_9 + \hbar\,P_9$ (Yang's $R$-matrix).
\item YBE verified numerically on $\mathbb{C}^{27}$ at four test points.
\item Level-$1$ RTT: $36$ independent relations ($= \mathfrak{gl}_3$ bracket),
 of which $21$ nontrivial and $15$ commuting pairs.
\item For $\mathfrak{sl}_3$: $28$ independent relations per order.
\item Triangle sectors: $4$ configurations
 ($T^L_{12}$, $T^L_{21}$, $T^R_{12}$, $T^R_{21}$).
\item Triangle coefficient: $\frac{1}{2} = B(1,2)
 = \int_0^1(1-t)\,dt$.
\item Serre relations: $2$ (one for each simple root pair).
\item The RTT presentation equals the $d^2 = 0$ condition on
 $\Barch_3^{\mathrm{ord}}(\widehat{\mathfrak{sl}}_{3,k})$,
 confirming Construction~12.1.1 of the manuscript.
\end{enumerate}
